\documentclass[../main.tex]{subfiles}
\begin{document}
%==========================================TIÊU ĐỀ TẬP========================================
\begin{table}[h]
  \begin{adjustwidth}{-1cm}{}
    \begin{tabular}{>{\centering\arraybackslash}p{6cm}|>{\raggedright\arraybackslash}p{12.5cm}}
      \multirow{5}{6cm}
      % Cột bên trái: ÉPISODE và số tập
      {\\[-40pt]
        \flaregothic\fontsize{20pt}{20pt}\selectfont \centering
        {\contour{errortwo}{\textcolor{black}{ÉPISODE}}}\color{black}\\
        \burbank\fontsize{72pt}{72pt}\selectfont
        \includegraphics[height=3cm]{../../../../Image/Book?/number/num2.png}
        \\[18pt]
        % \flaregothic\fontsize{14pt}{14pt}\selectfont
        % \color{epattr}BIENVENUE, \uline{\color{Banpire} Mel} !
        % \flaregothic\fontsize{18pt}{18pt}\selectfont
        % \color{epattr}ÉPISODE PERDU
        \color{black}
      }
       &
      % Dòng thứ nhất: tên tập tiếng Nhật
      {\fotskip\fontsize{18pt}{18pt}\selectfont \ruby{転}{てん}\ruby{校}{こう}\ruby{生}{せい}の\ruby{噂}{うわさ}
      }                       \\[6pt]
      % Dòng thứ hai: tên tập tiếng Anh
       & {\cafeteria\fontsize{18pt}{18pt}\selectfont The Rumor About Transfer Students}                                 \\[4pt]
      % Dòng thứ ba: tên tập tiếng Việt
       & {\vnmsans\fontsize{18pt}{18pt}\selectfont Lời đồn về học sinh chuyển trường}                                   \\[8pt]
      % Dòng thứ tư: Nhân vật xuất hiện
       & {\caviar{\fontsize{14pt}{14pt}\selectfont Track used: The Desert (High Flyer - Dave James \& Keith Beauvais)}}
      \\[18pt]
      % Dòng cuối cùng: Thời gian diễn ra của tập (thường sẽ là ngày phát hành)
       & {\continuum\fontsize{16pt}{16pt}\selectfont Ngày phát hành: 26/02/2022}                                        \\[8pt]
    \end{tabular}
  \end{adjustwidth}
\end{table}
%===========================================NỘI DUNG==========================================
\color{black}

\pov{Haragen Saikyou}

Suốt cả ngày hôm đó, cứ mỗi giờ ra chơi hay nghỉ giữa tiết, lại có thêm vài gương mặt mới tiến đến chỗ tôi và Nii-chan. Có người chỉ tò mò hỏi han, có người chủ động bắt chuyện, thậm chí còn nói thẳng là muốn kết bạn.

Thực lòng mà nói, tôi chưa từng nghĩ bản thân sẽ trở thành tâm điểm chú ý như thế này. Cảm giác đó giống như bị ném thẳng vào vị trí của một người nổi tiếng hay một KOL nào đó, lúc nào cũng có người muốn bắt chuyện hay chụp hình chung. Nhưng khác với sự chú ý kiểu thị phi, đây lại là những lời chào mừng và nụ cười thật sự, ít nhất thì, tôi muốn tin là vậy.

Ngày học cứ trôi qua trong vòng lặp đều đặn: học bài, gặp gỡ thêm vài người mới, rồi lại quay về với việc học. Mọi thứ diễn ra nhanh đến mức tôi gần như không kịp cảm nhận rõ ràng, chỉ biết là mỗi lần chuông báo tiết mới vang lên, tôi lại có thêm một hoặc hai người để nhớ mặt gọi tên.

Và rồi, chuông báo hết giờ học reo lên, kết thúc ngày đến trường đầu tiên của hai anh em chúng tôi.

Nii-chan ở lại lớp, vì Saki-san và Hanabi-san vẫn chưa chịu bỏ qua cơ hội ``khai thác thông tin'' để chuẩn bị cho bữa tiệc chào mừng vào cuối tuần sau.

``Vậy em về trước nha, Nii-chan!'' tôi nói, khoác cặp lên vai.

``Ừ. Nhớ nấu cơm sớm vào đấy,'' anh cười đáp, giơ tay vẫy nhẹ trước khi quay lại nói chuyện tiếp với hai cô bạn mới quen.

Khi bước ra cửa, tôi vẫn nghe loáng thoáng tiếng trò chuyện của họ, âm thanh hòa lẫn với tiếng bàn ghế dịch chuyển và tiếng giày bước khắp lớp. Có vẻ cuộc bàn luận đó sẽ còn kéo dài một lúc nữa.

Tiến ra hành lang, tôi bắt gặp hai cô gái mà chúng tôi cũng vừa quen không lâu, Inari-san, với bờm tai cáo trắng dễ thương nhưng lúc nào cũng hơi rụt rè, và Touka-san, nàng tóc xanh tràn đầy năng lượng.

``Saikyou-san! Muốn về với bọn mình không?'' Touka-san gọi lớn, giọng đầy hào hứng.

``M-Mình cũng muốn đi về với cậu!'' Inari-san cũng lên tiếng, dù giọng vẫn mang chút ngại ngùng.

Tôi mỉm cười, đáp ngay: ``Đi, đi chứ!''

Thế là ba chúng tôi rời khỏi trường cùng nhau. Chẳng biết bắt đầu từ lúc nào, câu chuyện đã chuyển sang đủ thứ đề tài, từ bài tập hôm nay, mấy chuyện vặt trong lớp, cho đến những tin đồn linh tinh mà học sinh ở đây truyền miệng. Tiếng cười, tiếng trêu chọc, và cả những câu than thở về việc làm bài tập như thể đang dệt thành một khúc nhạc quen thuộc của tuổi học trò.

Có lẽ, đây chính là điều mà người ta hay gọi là ``khoảnh khắc đời thường'' của học sinh nữ, một nhóm bạn cùng nhau tán gẫu trên đường về, tay khua khoắng theo từng câu chuyện, tiếng cười vang lên giữa con phố nhỏ.

Và tôi chợt nhận ra\dots\ hình như đây là lần đầu tiên tôi trải qua cảm giác này. Lần đầu tiên được hòa mình vào một cuộc sống nữ sinh bình thường đến vậy.

\ct

Kể cả khi đã bước ra khỏi cổng trường, những câu chuyện rì rầm của ba đứa con gái chúng tôi vẫn chưa có dấu hiệu chấm dứt. Thực ra thì, cũng chỉ là giữa tôi với Touka-san thôi, còn Inari-san vẫn im lặng nhiều hơn, thỉnh thoảng mới chen một câu, ánh mắt như đang theo đuổi một suy nghĩ nào đó.

Chúng tôi cứ thế nói đủ thứ tạp nham, từ bài tập về nhà cho đến mấy tiệm ăn ngon trong thị trấn. Con đường về nhà được nhuộm một màu vàng cam ấm áp của hoàng hôn, ánh sáng xuyên qua những hàng cây hai bên đường, tạo thành những vệt sáng dài như đang dẫn lối.

Hôm nay đúng là một ngày đi học khác hẳn so với những lần trước. Không phải kiểu cắm đầu học và về nhà ngay, mà là có người trò chuyện cùng, thậm chí còn được thong thả ngắm cảnh. Cảm giác này\dots\ thật mới mẻ, nhưng cũng đầy thích thú. Đã từ lâu rồi mới có thể trò chuyện với những người khác mà tôi có thể tin tưởng\dots

\dots dù không hẳn là hoàn toàn.

Có lẽ tôi đã bắt đầu hòa nhập với thị trấn nhỏ bé này. Sự yên bình nơi đây khác hẳn với nhịp sống vội vàng ở thành phố, một thứ quý giá mà thời buổi này khó mà tìm được.

``Cậu đã thật sự thấy thoải mái ở đây chưa?'' Touka-san bất chợt ngoái lại, nụ cười của cô ấy nhẹ như gió chiều.

``Tôi nghĩ là rồi đấy. Dần dà cũng bắt đầu thấy hứng thú,'' tôi đáp, lòng chợt nảy ra một danh sách những nơi muốn ghé thăm. Nhưng vẫn còn nhiều chỗ ở thị trấn Aogami này tôi chưa từng đặt chân đến.

``A-A-nô\dots'' Inari-san đột ngột lên tiếng, giọng run nhẹ như đang dò từng từ một. Dường như cô ấy muốn nhập cuộc nhưng lại chưa biết mở lời thế nào. Sau vài giây ấp úng, cuối cùng cô cũng nói: ``Cậu đã từng đến Công viên Aogami số 2 chưa?''

Tôi hơi ngẩn ra. ``Chưa. Sao thế?''

Khuôn mặt nàng ``Cáo trắng'' giãn ra, như vừa thở phào. ``Hồi nhỏ mình với Hanabi-chan hay chơi ở đó,'' cô nói, lần này liền mạch hơn. ``Chỗ đó rộng lắm, cảnh đẹp nữa\dots\ Nếu rảnh thì cậu thử đến nha!''

Tôi chưa bao giờ thấy công viên ở thị trấn thì thế nào, nên cũng khá tò mò. ``Hôm nào rảnh tôi sẽ rủ Nii-chan với các cậu cùng đi!''

Thật sự\dots\ hai người này thân thiện đến mức tôi không thể không mỉm cười. Điều ước ``có một người bạn'' của tôi và Nii-chan, xem ra, đang thành sự thật.

Chúng tôi đi thêm một đoạn thì đến bên bờ sông cắt ngang thị trấn. Nước sông trong vắt đến nỗi có thể thấy vài con cá lấp ló dưới đáy, bơi chậm rãi theo dòng. Ở phía bên kia là một khu rừng xanh mướt, tán lá cao thẳng tắp, mang một vẻ đẹp yên tĩnh đến bất ngờ.

\img{e1-2a}

``Đẹp thật đó\~{}'' Inari-san khẽ thốt lên, đôi mắt phản chiếu mặt nước lấp lánh. Ở đây, không có tiếng còi xe inh ỏi, chỉ có tiếng nước chảy và gió lùa qua kẽ lá.

``Mình còn thấy mấy bụi hoa đang nở ở bờ bên kia kìa,'' Touka-san chỉ tay, nơi vài chùm hoa màu trắng và vàng đang rung rinh giữa tán lá rậm rạp.

``\dots Aogami quả là một nơi biết cách khiến người ta bất ngờ,'' tôi gật đầu. Trong cái tĩnh lặng này, vẫn ẩn giấu một điều gì đó khó gọi tên, nhưng khiến tim tôi rung lên nhè nhẹ.

Tôi thấy ở góc bờ sông có một lối mòn nhỏ, lẩn khuất giữa những bụi cây, dẫn thẳng vào rừng. Tôi chỉ tay vào chúng, rồi hiếu kỳ hỏi: ``Còn đường kia là đi đâu vậy?''

Cô bạn tóc xanh đáp: ``Nếu đi theo nó rồi vòng sang bờ bên kia, cậu sẽ tới ngôi đền Aogami.''

``Ngôi đền Aogami? Ở đây cũng có đền à?''

``Ừ! Cũng là một nơi đáng để ghé lắm,'' Inari-san nói, giọng đã bớt ngập ngừng.

``Vậy nó đặc biệt ở điểm nào?'' tôi hỏi tiếp.

``Ờm\dots\ Khó giải thích lắm,'' cô bạn Thần đồng gãi đầu. ``Tớ cũng không có nhiều kỷ niệm về đó\dots''

``Với lại, bọn mình cũng chẳng rõ nó được xây từ bao giờ. Chỉ biết là nó ở đó rất lâu rồi,'' Lớp trưởng thêm vào.

Tôi định hỏi thêm, nhưng cả hai đều có vẻ không muốn nói sâu. Tôi chỉ nhún vai: ``Thôi, vậy không hỏi nữa.''

\dots

\dots

`\dots Cùng nhau là bạn nhé\dots''

``\dots Làm bạn với họ đi\dots''

Một thoáng gì đó lướt qua tai tôi. Gió? Lá xào xạc? Hay là\dots\ một giọng nói nào đó?

Cùng lúc đó, từ lối nhỏ kia, một luồng khí lạnh mơ hồ như len vào da thịt. Giống như có thứ gì đó đang âm thầm quan sát.

Tôi chớp mắt. Rõ ràng nơi này ban nãy còn đẹp như tranh, vậy mà giờ lại thấp thoáng một bóng tối khó tả.

Vả lại, cả Inari-san lẫn Touka-san đều không bảo tôi nên tới ngôi đền. Như thể họ đang giấu nhẹm một cái gì đó.

\dots

\dots

Chắc là do tôi tưởng tượng thôi. Vả lại, nếu họ không bảo, tôi cũng không nhất thiết phải bén mảng tới đó.

``V-Vậy\dots\ chúng ta về thôi nhỉ?'' tôi cất tiếng, cố nặn ra một nụ cười nhưng lại cảm thấy mồ hôi lạnh đang rịn ra sau gáy.

Hai cô gái khẽ nghiêng đầu nhìn tôi. Ánh mắt của họ không hẳn là nghi ngờ, mà giống như đang dò xét xem tôi có phát hiện ra điều gì không. Nhưng chẳng ai hỏi thêm. Có lẽ\dots\ họ cũng chẳng muốn tôi đào sâu.

``Ừ, về thôi,'' nàng tóc xanh khẽ gật đầu.

``Bọn mình cũng nên về. Saikyou-san, anh trai cậu vẫn còn ở trường phải không? Không biết Hanabi-chan chuẩn bị xong cho bữa tiệc chưa\dots''

``Mình còn sợ là cậu ấy không nhớ đường về nữa kìa---''

``Không cần bận tâm, không cần bận tâm!'' tôi vội xua tay, cười xòa như để xua tan bầu không khí lạ lùng đang bao trùm. ``Nii-chan thuộc đường rồi, ổn thôi mà.''

Cả hai cùng mỉm cười. ``May thật,'' Inari-san nói, nhưng giọng cô ấy lại giống như đang thở phào vì\dots\ lý do gì khác.

Không ai nói gì thêm, chúng tôi lặng lẽ bước tiếp.

\ct

Tôi biết chắc mình không phải người duy nhất cảm thấy có gì bất thường. Từ lúc rời bờ sông, hai cô bạn bắt đầu trò chuyện về những hiện tượng kỳ lạ trong thị trấn, nhưng không một lần nhìn thẳng vào tôi.

``Nhắc mới nhớ\dots\ hôm nay lớp bên cạnh vắng nhiều người lắm,'' Inari-san lên tiếng, nhưng cách cô ấy nói như đang thử dò xem phản ứng của tôi.

Touka-san thì đáp với vẻ bình thản đến khó chịu, như thể chuyện đó chẳng còn gì mới mẻ:
``Với cả\dots\ chẳng ai liên lạc được với họ nữa.''

``T-Thế à?'' tôi cố chen vào, nhưng họ như chẳng mấy bận tâm đến lời tôi.

Đến khi tôi nhắc lại chuyện lớp vắng, họ mới giật mình và cười gượng: ``P-Phải ha\dots\ hôm nay lớp mình cũng yên tĩnh hẳn\dots''

Kể ra thì cũng lạ. Suốt từ lúc tan trường, tôi gần như không gặp ai khác. Ngay cả cô bạn có bờm tai chó hồi sáng ở cổng trường cũng biến mất không dấu vết.

``Hai cậu\dots\ có nghĩ đó là cúm mùa không?'' nàng tóc trắng hỏi, nhưng mắt lại đảo sang hướng khác.

``M-Mình nghĩ vậy. À, Saikyou, nhớ giữ gìn sức khỏe nha,'' nàng Thần đồng nói, mồ hôi khẽ rịn trên thái dương.

``Ừ\dots\ Nhưng lạ thật. Cúm thì làm gì mà hàng loạt người nghỉ như thế được?''

``Tớ cũng thấy lạ,'' Inari đáp. ``Chỉ là cúm thôi mà\dots''

``Hẳn phải có lời giải thích hợp lý hơn chứ?'' tôi buột miệng nói ra suy nghĩ của mình.

Ngay khoảnh khắc đó, Touka dừng bước. Cô xoay người lại. Ánh mắt ấy\dots\ khiến tôi sởn da gà.

\img{e1-2b}

Một cái trừng mạnh đến mức như đóng băng không khí xung quanh. Không còn chút ấm áp nào của người bạn vừa trò chuyện vui vẻ ban nãy. Đôi mắt vàng ấy ánh lên thứ gì đó lạnh lẽo và sắc bén, đủ để khiến tôi quên mất mình đang thở.

``\emph{Chỉ là cúm mùa thôi},'' giọng cô ấy thấp, chậm, và tuyệt đối không để cho tôi phản bác.

Trong thoáng chốc, tôi thấy khóe môi cô giật nhẹ, như thể cảnh báo.

\dots

\dots

Trông họ đáng sợ thật. Cái cảm giác ấy\dots\ không phải là giận đơn thuần, mà là muốn dập tắt ngay từ gốc bất kỳ nghi ngờ nào của tôi. Cứ như thể họ phải làm thế.

Tôi không hiểu tại sao, nhưng trong khoảnh khắc đó, tôi biết chắc mình vừa chạm vào một điều mà họ không muốn ai nhắc đến.

Và họ sợ tôi biết sự thật.

\dots Hay sợ rằng tôi sẽ trở thành một phần của nó?

Tôi lặng người vài giây, rồi chỉ khẽ gật: ``Ờm\dots\ hiểu rồi. Nếu cậu cho là vậy.''

Bầu không khí ngột ngạt tan đi nhanh chóng, như thể chưa từng tồn tại. Họ quay lại nói cười với tôi, nhưng tôi có thể cảm nhận rõ, nụ cười ấy gượng gạo và có phần giả tạo đến mức nào.

Trong lòng tôi, một câu hỏi lởn vởn không chịu biến mất: Aogami\dots\ rốt cuộc nơi này đang che giấu cái gì?

Gió chiều thổi nhẹ qua con phố, cuốn theo tiếng lá xào xạc.

Tôi không hề biết rằng, ở khoảng cách chỉ vừa đủ để tránh khỏi tầm nhìn của chúng tôi, có ai đó đang đứng ở gần đó theo dõi mọi thứ.

Đôi mắt nâu lặng lẽ dõi theo, không một biểu cảm, nhưng lại mang một thứ ánh nhìn như thể đã chứng kiến mọi thứ, và đang cân nhắc điều gì đó.

Bóng cô hòa vào dãy nhà phía sau, mỏng manh như sương\dots\ nhưng tuyệt đối không thể coi là vô hại.

Có gì đó thật sự đáng ngờ. Hai cô bạn mà tôi vừa quen - Inari và Touka - hiển nhiên cũng nằm trong ``danh sách cần để mắt''. Nhưng thay vì chất vấn, tôi chọn im lặng, tiếp tục bước cạnh họ. Không khí buổi chiều ấm áp, nhưng cảm giác trong lòng tôi lại chẳng hề dễ chịu chút nào.

Bất chợt, Lớp trưởng khẽ gọi: ``Ê-tồ\dots\ Ờm\dots'' Giọng cô nhỏ xíu, như muốn nói gì nhưng lại thôi.

Tôi quay sang, định đáp, nhưng vẫn giữ im lặng, chỉ ``ừm'' một tiếng nhẹ.

Cô bạn Thiên tài nghiêng đầu, vẻ ngạc nhiên: ``Inari-san, Saikyou-san, sao thế?''

Inari-san vội xua tay: ``À, không, chỉ là\dots\ quả thật trường mình dạo này có nhiều người nghỉ học lắm nhỉ.''

Tôi gật nhẹ, như muốn tập trung vào chủ đề này: ``Chả trách mà lớp bên cạnh lại im ắng thế. Mà Touka-san nói rằng họ bị cảm, đúng chứ?''

Cô bạn tóc xanh trầm ngâm một lúc rồi đáp: ``Mình nghĩ vậy. Cơ mà nhá\dots\''

Inari-san chen vào, tiếp lời: ``Dạo này bọn mình còn chẳng nghe được tin gì từ những học sinh ở lớp bên cạnh nữa.''

Touka-san gật đầu, hơi nhíu mày: ``Đó cũng là vấn đề lớn đấy\dots\ Mình tin rằng chắc chỉ là do cơn cảm cúm mùa hè thôi.''

Lớp trưởng nghiêng người nhìn cô bạn: ``Thế á?''

``Ừ.'' cô trả lời ngắn gọn.

Tôi mím môi: ``Nhưng mà\dots\ tôi lại cảm thấy khó tin lắm. Tôi biết là cậu vừa trừng mắt với tôi\dots''

Cô bạn tóc xanh hơi khựng lại, mắt chớp nhanh khi tôi nhắc về chuyện ban nãy: ``\dots Mình xin lỗi. Mình không cố ý dọa cậu sợ.''

Tôi nhẹ giọng: ``Ừ, tôi biết. Nhưng mà, tại sao cậu lại nghĩ thế?''

Cô bạn Thần đồng thở ra, giải thích: ``Cậu biết mấy ngày gần đây trời càng nóng hơn chứ?''

Tôi suy nghĩ một thoáng: ``Giờ cậu nói tôi mới để ý\dots''

Lớp trưởng cũng gật gù: ``Mình cũng thấy hơi khó chịu thật.''

``Thì đó,'' Touka nói tiếp, ``trời nóng thì người ta sẽ ăn mất ngon hơn, ngủ cũng ít hơn nữa. Thế nên nhịp vận hành của cơ thể họ cũng sẽ bị ảnh hưởng theo.''

Tôi chậm rãi gật đầu: ``Cậu nói cũng có lý\dots''

Inari-san lắp bắp: ``Đ-Đúng nhỉ. Cơn cảm cúm mùa hè thôi nhỉ\dots\ Nếu như vậy thì may quá chừng\dots''

Dù vậy, Touka-san vẫn giữ giọng nghiêm túc: ``Mình thậm chí còn chẳng biết liệu chúng ta có nên an tâm về chúng không đây. Nghe đồn trần cúm này diễn ra dai dẳng lắm.''

Tôi khẽ thở dài: ``Tôi lo cho Nii-chan quá. Anh ấy có thể chất yếu mà.''

Lớp trưởng tròn mắt ngạc nhiên: ``T-Thế ư? Trông cậu ấy cũng có đâu đến nỗi\dots''

``Ừ,'' tôi đáp, ``trông anh ấy khỏe thế thôi chứ anh ấy đã ốm là không ngóc đầu dậy nổi luôn. Tôi chăm anh ấy nhiều, tôi biết chứ.''

Touka-san mỉm cười hơi chua chát, mắt ngước lên trời như hồi tưởng lại: ``Hồi năm ngoái mình cũng thức đêm học nhiều quá, thế là thi xong mình nằm ốm liệt giường ba ngày liền. Mình tưởng mình không qua môn được cơ.''

``Giống như Nii-chan hồi thi đầu vào nhỉ,'' tôi bật cười.

Inari-san cũng góp lời: ``Chắc lúc đó cực lắm. Không thể đi đâu và làm gì được kia mà.''

``Đấy,'' tôi gật đầu, ``thế nên là mọi người giữ gìn sức khỏe nha.''

``Ừ-Ừm,'' Lớp trưởng đáp, rồi bỗng nhớ ra điều gì, ``mà nhắc đến bài kiểm tra thì\dots''

Cô bạn tóc xanh cười nhẹ: ``Sắp tới lúc thi cuối kỳ rồi nhỉ? À, Saikyou-san với Kyuen-san mới vào nên chắc cậu không biết đâu.''

Tôi lắc đầu: ``À không, cái này thì hai anh em bọn mình đã chuẩn bị từ trước rồi. Các cậu không phải lo.''

``May thật đó,'' Touka-san nói với sự nhẹ nhõm, rồi quay sang Inari, ``Inari-san, cậu lại định thức đêm học bài nữa sao?''

``Chắc là như vậy quá\dots'' cô bạn tóc trắng gãi má.

Cô bạn Thần đồng nhíu mày: ``Mình biết là cậu có nhiều thứ để nhồi nhét vào đầu lắm, nhưng mà chú ý nghỉ ngơi đầy đủ.''

``Ừ\dots\ Có lẽ vậy,'' Inari thở ra.

Tôi cười khẽ: ``Cậu là lớp trưởng mà, làm gương cho người ta đi chứ.''

``Như Saikyou-san nói đó,''Touka-san cũng đồng tình. ``Không có cậu thì tới mỗi tiết học, cả lớp sẽ thành cái chợ mất.''

Tôi bật cười: ``Giống như mấy lớp học trước mà chúng tôi đã trải qua nhỉ\dots\ Mà, Trưởng ơi, nhớ giữ gìn sức khỏe ha? Có gì khó thì cứ hỏi hai anh em bọn tôi!''

Inari cười trừ, khoát tay: ``Cảm ơn hai cậu. Được rồi, mình sẽ cố gắng chăm sóc bản thân tốt hơn.''

Tôi giơ ngón tay chỉ thẳng: ``Hứa đấy nhá!''

``Mình hứa, mình hứa,'' Inari bật cười.

\ct

Đến ngã rẽ, chúng tôi tạm biệt nhau. Touka-san và Inari-san cùng nhau bước tiếp về phía khu nhà phía đông, còn tôi quay về hướng ngược lại. Tiếng trò chuyện rì rầm của hai người họ mờ dần trong gió chiều.

Bước đi một mình, tôi lại cảm nhận rõ rệt sự bất an đang len lỏi. Cúm mùa hè\dots\ nghe thì hợp lý, nhưng cái cách họ nhìn nhau trước khi trả lời, hay khoảnh khắc Touka trừng mắt với tôi, tất cả đều chẳng giống chuyện ``cúm'' thông thường.

Tôi siết chặt quai cặp, thở ra thật chậm.
``\textit{Có lẽ\dots\ mình sẽ phải để mắt đến chuyện này nhiều hơn nữa.}''

Bất chợt, ánh mắt tôi chạm vào một thứ lấp lánh bên vệ đường. Tôi dừng lại, cúi xuống.

Đó là một mảnh dây đeo đồng hồ, loại dây cao su màu đen nhám, ở phần đuôi có gắn một khung kim loại nhỏ, kèm theo một cái kim mảnh như có thể điều chỉnh để ôm sát cổ tay.

Tôi cầm nó lên, lật qua lật lại trong tay. Dây đã hơi xước, nhưng còn nguyên hình dạng.

\dots Lạ thật. Trông nó rất giống cái mà anh trai tôi đã kể - thứ anh nhặt được\dots\ trong mơ. Khác ở chỗ, cái dây đeo mà anh nhặt được gồm có lỗ, vậy tôi nghĩ dây này sẽ dùng để khít với cánh tay.

Gió chiều lùa qua, mang theo một thoáng lạnh buốt dọc sống lưng.

Tôi chợt thấy mình đang cầm một thứ\dots\ không nên xuất hiện ở đây.

Tôi cất nó vào túi áo, bước đi nhanh hơn.

Không biết vì sao, nhưng tôi có cảm giác rằng\dots\ nó có thể giúp chúng tôi giải đáp một cái gì đó.

\begin{center}
  \uline{Ngày hôm sau.\\Lớp 1-C, Trường Trung học Aogami.}
\end{center}

Ngày thứ hai chúng tôi tới ngôi trường này. Và thật sự, cũng chẳng dễ chịu hơn hôm qua là bao.

Không phải chỉ vì những hiện tượng kỳ bí hay mấy chuyện bất thường đã bắt đầu lởn vởn trong đầu tôi từ hôm qua.

Mà hôm nay còn mưa như trút nữa. Cũng may là hai anh em kịp chạy đến trường trước khi bị ướt nhẹp\dots\ chỉ có điều, lại quên không mang ô.

\img{e1-2c}

``Anh tưởng dự báo bảo hôm nay không mưa cơ mà?'' Nii-chan chống tay lên bàn, liếc qua cửa sổ đầy những vệt nước mưa.

``Đài dự báo láo đấy\dots'' tôi tặc lưỡi. Lúc sáng nhìn cái đám mây đen đặc quánh kia là biết ngay có chuyện rồi. Nhưng anh tôi cứ khăng khăng bảo ``chắc gì hôm nay đã mưa!'' và\dots\ như các bạn thấy rồi đấy.

Hôm nay cũng là một ngày kỳ quặc không kém. Ngoài việc hôm nay tiếp tục là những tiết học không có gì đặc biệt, giống như hồi ở các trường mà hai anh em chúng tôi đã theo học trước đó\dots

``Mọi người\dots\ hôm nay nghỉ hết à?''

\dots lời anh tôi nói vô thức bật ra lại đúng như những gì hiện tại. Chỉ có hai anh em tôi, Hanabi-san và Suzune-san đang ríu rít ở góc lớp, Inari-san và Touka-san thì thì thầm với nhau, còn Saki-san thì ngồi im lìm một mình.

``Em cũng nghĩ thế\dots'' tôi đáp nhỏ.

``Chắc lại ốm vặt gì thôi?''

``C-Chắc vậy?'' tôi cười gượng, nhưng cổ họng hơi nghẹn lại.

Cảnh tượng này\dots\ sao quen quá. Giống như tôi đã thấy ở đâu đó\dots\ trong mơ. Cũng là những con người này, dù rằng chỉ khác vị trí ngồi và tư thế. Và trực giác mách rằng, đây không phải là trùng hợp.

Trong đầu tôi vẫn vang lại khoảnh khắc hôm qua. Tôi vẫn không quên được ánh nhìn sắc lạnh bất ngờ của Touka-san, đi kèm câu nói khẽ buông ra: ``Chỉ là một trận cúm mùa thôi.''

Cái cách mà cô ấy hơi khựng lại, còn Inari-san thì liếc sang với ánh mắt lúng túng khi tôi hỏi vặn\dots\ cả hai người họ rõ ràng không thoải mái khi nhắc đến chuyện đó.

Không, không thể chỉ là một trận cúm thường. Nếu là cúm mùa thật, đâu đến mức cả lớp tôi và lớp bên cạnh đồng loạt biến mất như thế. Dù sao\dots\ họ cũng quả quyết vậy, nên tạm thời tôi đành giả vờ tin.

``\dots Saikyou.''

Tiếng gọi cộc lốc làm tôi giật mình. Anh trai tôi đang nhìn chằm chằm, đầu hơi nghiêng một bên như đang đọc thấu suy nghĩ tôi.

``Có chuyện gì mà mặt em đần thối ra thế?''

``K-Không có gì!''

``Hừm\dots\ nghi lắm nha\dots'' anh nghiêng người áp sát, đôi mắt săm soi từng chút. ``Có gì đó không ổn đúng không? Nhìn mặt em là biết.''

\dots Tôi vốn không thể giấu nổi anh song sinh. Chỉ cần nhìn là anh biết ngay.

``Anh cũng thấy lạ giống em thôi,'' Nii-chan nói, lùi ra. ``Nói thì hơi đa nghi, nhưng việc cả lớp nghỉ thế này\dots\ thật sự không bình thường đâu.''

Đúng là tinh ý. Đến nước này, tôi cũng chẳng muốn giấu nữa. Lớp học này, thị trấn này, rồi mọi thứ xung quanh dường như đang hòa vào một bản nhạc kì quái mà tôi chưa nghe bao giờ.

``Em thấy\dots\ kỳ lắm. Chúng ta còn chẳng thấy ai ngoài sáu người bọn mình.''

``Chưa kể Saki-san kia\dots'' Nii-chan liếc về cô nàng tóc vàng hai bím đang khoanh tay bất động ở góc lớp, trông như đang ngủ gật.

``Thậm chí các thầy cô cũng chẳng thấy gì lạ. Ánh mắt họ\dots\ vô hồn hẳn.'' Tôi rùng mình khi nhớ lại.

Bạn học của tôi hành xử khác hẳn những người bình thường. Như thể\dots\ họ đang bị điều khiển.

\dots Tôi có nên hỏi thẳng họ không? Một nửa tôi muốn, một nửa thì không. Ánh nhìn hôm qua của Touka-san\dots\ vẫn khiến tôi chùn lại.

\dots

Đang còn chìm trong mớ suy nghĩ hỗn độn, một giọng nói rộn ràng bất ngờ vang lên từ phía sau, kéo tôi và Kyuen-nii-chan trở về thực tại.

``Ê, ê này, hai cậu ơi!''

Quay lại, tôi thấy Hanabi-san đang vẫy tay gọi, nụ cười sáng bừng như xé tan bầu không khí mưa u ám ngoài kia. Bên cạnh cô là Suzune-san, đang nghiêng người tò mò. Không hiểu họ định kéo chúng tôi vào chuyện gì nữa đây.

``Mình cũng vừa bàn chuyện này với Suzune-san rồi,'' Hanabi-san nói, giọng đầy ẩn ý. ``Không biết các cậu sẽ phản ứng thế nào\dots''

``Rồi, rồi. Nó là cái gì vậy?'' tôi nhướng mày, cảm giác như sắp bị lôi vào một trò gì đó.

``Trông mặt cậu hớn hở thế kia chắc là có cái gì vui lắm hả?'' Nii-chan lên tiếng, giọng vừa trêu vừa dò xét.

Suzune-san khúc khích cười, rồi bất ngờ đưa tay quạt quạt trước mặt. ``Chuyện là thế này\dots''

Tôi chợt nhận ra rằng\dots\ quả thật, dù trời đang mưa nhưng hơi nóng vẫn quẩn quanh khắp lớp học. Không khí bí bức đến khó chịu. Gần đầu hè rồi, bảo sao\dots

``Nóng quá\~{}!'' Suzune-san rên rỉ, rồi buông tay xuống bàn như thể sắp tan chảy.

``Công nhận đúng là trời nóng thật đấy,'' anh tôi đồng tình, khẽ kéo cổ áo. ``Trời mưa tầm tã ngoài kia mà chẳng thấy mát tí nào\dots''

``Ừ thì\dots\ Cũng sắp vào hè rồi mà,'' Hanabi-san gật gù.

``Trời nóng thế này học đúng là khó vào thật á\dots'' cô bạn Nhiệt thành than tiếp, giọng dài lê thê.

``Thật. Ở đây không có điều hòa hay thậm chí là quạt trần à\dots?'' tôi thở dài, rồi liếc sang Suzune. ``Mà này, Suzune-san, tôi nghe nói vốn dĩ cậu lười học vãi cả linh hồn ra, còn kêu ca gì.''

``Ể!?'' Suzune bật dậy như bị điện giật, mắt tròn xoe.

``Em ấy nói đúng mà,'' anh tôi hùa vào, miệng nhếch lên. ``Có quạt hay không có quạt thì cậu vẫn lười như nhau thôi.''

Suzune phồng má, đôi mắt long lanh như muốn phản pháo nhưng không tìm ra lý lẽ. ``Hai người đúng là xấu tính\dots\ Đến khi nào ta mới được nghỉ hè đây?''

``Thi xong cuối kỳ thì mới nghỉ được chứ,'' anh tôi đáp tỉnh bơ.

``Thì đó,`` Hanabi-san tiếp lời, rồi ánh mắt bỗng sáng lên. ``Mà nhắc đến nghỉ hè thì\dots\ các cậu đã có dự định gì chưa?``

``\dots Chắc là tự nhốt trong nhà chơi điện tử hưởng gió điều hòa thôi nhỉ?'' Nii-chan nói như thể đó là chân lý bất di bất dịch.

``Mình lại muốn đi ra ngoài chơi hơn,'' tôi chen vào, rồi hướng sang cô bạn năng nổ. ``Hanabi-san chắc là ưa các món hoạt động ngoài trời hơn nhỉ?''

``Ưm!'' vô nàng gật mạnh, đôi mắt long lanh. ``Cắm trại trên núi, rồi được đi bơi ở hồ\dots\ mình có nhiều thứ muốn làm khi nghỉ hè lắm chứ!''

``Thấy chưa nào, Nii-chan!'' Tôi lập tức quay sang trêu anh. ``Ra ngoài trời thì mới có thú vui để chơi chứ!''

``Rồi rồi, khỏi nói. Anh cũng đếch quan tâm đâu,'' Kyuen gạt phắt, chẳng mảy may bị dao động.

``Kyuen-san đúng là lười vận động ha\~{}'' cô bạn năng nổ chun mũi trêu lại.

``Sao cũng được. Suzune-san thì sao?''

Vừa được hỏi tới, Suzune-san liền chống tay lên bàn, đôi mắt sáng rực: ``Với mình thì mình mong chờ lễ hội nhất! Mình mong chờ ngày hôm đó quá đi!''

``Lễ hội\dots? Kiểu, hội hè á?'' tôi nghiêng đầu.

``Đương nhiên rồi!'' Hanabi-san đáp nhanh, giọng đầy hào hứng. ``Cứ mỗi năm thì thị trấn Aogami sẽ tổ chức nó một lần. Các cậu đến đúng thời điểm lắm đấy!''

``Trông các cậu hớn hở vậy, chắc là họ chuẩn bị chu đáo lắm nhỉ?'' anh tôi ngóc đầu hỏi, giọng nửa tò mò nửa châm chọc.

``Tất nhiên!'' Hanabi-san gật đầy quả quyết. ``Mình nghe nói năm nay mọi người trong thị trấn nỗ lực chuẩn bị lắm đấy.''

``Thế á? Tại sao vậy?'' -san nghiêng đầu. ``Năm nay là năm đặc biệt à? Kiểu, kỷ niệm 50 năm thành lập lễ hội này sao?''

``Hmmm\dots\ Tôi cũng không biết nữa,'' cô bạn năng nổ chớp mắt, rồi phẩy tay. ``Chắc chả để ý.''

``Này. Tôi tưởng là lễ hội nào họ cũng chuẩn bị kỹ lưỡng lắm chứ? Có gì ngạc nhiên đâu?'' tôi hỏi lại.

``Không phải thế,'' Hanabi-san nghiêng người về phía chúng tôi, giọng hạ xuống một chút như sắp kể bí mật. ``Mình nghe đồn năm nay họ làm khác hẳn so với những năm trước! Chắc là nó có liên quan gì đó đến Thần họ thờ phụng chăng?''

``Ồ\dots\ Nghe có vẻ xoắn não đấy,'' Suzune-san xuýt xoa.

``Anh em tôi không tin vào Thần linh lắm đâu,'' tôi khoanh tay. ``Cùng lắm chắc là ngày kỷ niệm đặc biệt gì đó thôi.''

``\dots Cậu nói cũng phải. Thôi, quên những gì mình hỏi đi,'' Suzune gãi má, dường như hoàn toàn theo những gì mà tôi nói.

``Hể\~{} Các cậu có chắc không vậy?'' cô bạn năng nổ hơi nhướng mày. ``Không tôn trọng Thần linh kẻo sau này bị nghiệp quật đấy.''

``Cụ thể là như thế nào?'' anh tôi hỏi, nửa đùa nửa thật.

``Kiểu như là\~{} bắt các cậu làm bài tập hè gấp đôi chẳng hạn!''

``\dots Nghe xàm xàm thế nào ấy,'' tôi lầm bầm. Nii-chan cũng chỉ lắc đầu ngán ngẩm. ``Trò đó chỉ dọa được mấy đứa trẻ con lười làm bài tập  thôi.''

``Eek! Sợ thế!'' Suzune-san thì\dots\ rụt cổ, rồi giơ cả hai tay lên trời, lắc qua lắc lại như đang cầu khấn. ``Tôi hứa là sẽ thờ phụng ông mà! Tôi sẽ tham gia vào lễ hội, nên là đừng trừng phạt tôi mà\~{}!''

Cả hai chúng tôi đổ mồ hôi hột, cau mày với phản ứng có phần ngây thơ tới ngớ ngẩn: ``Cậu thực sự tin sái cổ lời cậu ấy nói sao, Suzune-san!?''

``Suzune-san đúng là ngố tàu mà,'' Hanabi-san che miệng cười.

``Gì chứ!?'' cô bạn Nhiệt thành đỏ mặt phản đối, như thể bị đem ra làm trò cười.

``Vậy chốt lại ý các cậu là sao đây?'' tôi chen ngang, muốn kết thúc màn trêu chọc này.

Hanabi-san lập tức xoay hẳn người về phía anh em tôi, phớt lờ luôn cô bạn tóc nâu đang vò đầu bứt tai. Cô dang tay, nụ cười tươi rói: ``Mình muốn rủ hai cậu đi cùng!''

Tôi chống cằm, suy nghĩ một lúc. Với một đứa thích ra ngoài như tôi, ý tưởng này thật sự hấp dẫn. ``Dĩ nhiên rồi.''

Rồi tôi liếc sang ông anh đang nằm dài như xác sống: ``Có khi Nii-chan ngốc nghếch của tôi cũng chẳng đi được đâu. Anh ấy lười chảy thây thế kia mà.''

``Em nói chỉ có chuẩn,'' anh đáp lại tỉnh queo, như thể đó là sự thật hiển nhiên, nên tôi không ngạc nhiên mấty. Nhưng hai cô bạn thì lại sững người trước câu trả lời thẳng tuột kia, cùng nhau đồng thanh: ``Thẳng thắn luôn kìa!''

Tôi cười xòa, tỏ vẻ cảm thông cho họ: ``Thôi thì kệ anh ấy đi. Thế lễ hội ở đây gồm có những gì?''

Hanabi-san chỉ mỉm cười, đứng dậy bước lại gần bàn hai anh em tôi. ``Thì\dots\ cũng giống mấy lễ hội khác thôi, có đủ quầy hàng các kiểu,'' cô đáp, giọng vừa hào hứng vừa giấu giếm điều gì đó.

Tôi rướn mày: ``Thế thôi sao? Không có gì đặc biệt hơn à?''

Cô nàng đưa mắt liếc sang hai bên, như thể đang tìm cách lảng đi. ``À\dots\ có thể là còn một cái gì đó, nhưng mà mình không chú ý lắm\dots''

Chuẩn bài rồi.

``Nó chắc cũng giống với mấy lễ hội mùa hè mà chúng ta từng đi thôi,'' Nii-chan lười biếng đáp, đầu vẫn gục xuống bàn.

Cô bạn Năng nổ liền lắc đầu: ``Không đâu. Đảm bảo sẽ khác hẳn! Với cả\dots\ đi chung mới vui mà!''

Tôi gật đầu: ``Cậu đã mời thì mình đi thôi. Nhưng Nii-chan á? Khó lắm.'' Tôi liếc mắt nhìn ông anh trai vừa lười ra ngoài vừa hướng nội của tôi, hỏi vọng lại: ``Nii-chan, có thật là anh sẽ không đi hử?''

``Anh miễn,'' anh tôi nói ngay, không thèm ngẩng mặt. ``Em cứ đi, khỏi lo cho anh.''

Nhưng có vẻ như Suzune-san và Hanabi-san lại không nghĩ vậy. Đặc biệt là cô nàng Nhiệt huyết, cô nhăn mặt bĩu môi: ``Bộ hai người quên Suzune này rồi sao? Tôi cũng muốn đi nữa!''

Cô lập tức cười tươi, sà tới như muốn kéo người ta khỏi chiếc ghế, rồi\dots

*BỐP!*

Suzune-san vung tay đấm nhẹ vào lưng Nii-chan. Nghe thì nhẹ, nhưng tiếng vang dội lại rõ mồn một. ``Dậy đi nào! Đừng có chây ì như thế!''

Tôi thoáng nghĩ đó là hành động vô hại, chỉ để trêu chọc, họa chăng cũng chỉ là thúc giục. Nhưng sắc mặt Nii-chan khẽ nhăn lại, một tay chống bàn, một tay đưa ra sau xoa xoa chỗ vừa bị đập.

``Ái\dots\ Nhẹ thôi chứ. Chỗ đó đau đấy.'' Dù cau có, anh vẫn thở dài, xuống nước nhượng bộ: ``Rồi, rồi, anh sẽ đi. Cho xong chuyện, được chưa?''

Suzune-san nghe anh ấy được bật đèn xanh, liền phấn khích nhảy tại chỗ: ``Anh em nhà Haragen sẽ đi chung với chúng ta!''

Tôi bật cười, và cũng mừng cho Nii-chan vì cuối cùng bị hai cô nàng ấy khuất phục\dots

Nhưng rồi\dots

\img{e1-2d}

Ánh mắt của Saki-san và Touka-san lập tức thu hút sự chú ý. Hai cô đứng như chôn chân, tròng mắt mở to, môi run run. Cô nàng chủ bán hoa thậm chí còn đưa tay che miệng, như thể vừa chứng kiến một điều cấm kỵ.

Cả căn phòng bỗng chốc hụt mất một nhịp.

``Ờ\dots\ Hai người sao thế?'' Nii-chan hỏi, nhìn bản mặt kinh hoàng mà mọi người đang thể hiện.

Ngay cả Suzune-san cũng chớp mắt khó hiểu: ``Gì vậy? Tôi làm gì sai à?''

Hanabi-san thì chỉ khẽ cau mày, im lặng.

Touka-san tiến lại một bước, giọng nghiêm hẳn, mặt tái mét đi: ``Suzune-san\dots\ cậu vừa làm gì thế?''

``Ể\dots? Thì\dots\ đấm nhẹ thôi mà? Có sao đâu\dots'' Suzune-san chớp mắt, rõ ràng không hiểu chuyện gì đang diễn ra.

``Không phải vậy!'' cô ấy liền cắt lời, giọng căng như dây đàn. ``Cậu thấy không? Kyuen-san đã cau mày rồi đó! Sao lại làm thế với học sinh mới!?''

``Ơ\dots\ M-Mình\dots\ ờm\dots'' cô bạn vừa vô thức đấm vào lưng anh tôi liền tỏ ra ấp úng.

Nii-chan thấy vậy, lập tức giơ tay xua: ``Không sao thật mà. Mình chỉ hơi giật mình thôi.''

``Đúng rồi, chỉ là cú đấm chơi thôi! Anh ấy không bị gì đâu,'' tôi vội thêm vào, cố làm dịu bầu không khí đang đặc quánh lại.

Nhưng chẳng ai cười nữa. Hanabi-san lén nhìn Inari-san, và tôi chợt nhận ra\dots\ cô nàng tóc trắng đang run nhẹ, mắt dán vào chỗ anh tôi vừa bị đánh. Suzune-san thậm chí còn cuống quýt nhìn về phía cô nàng tóc xanh đang bừng bừng sát khí rồi quay lại nhìn về phía chúng tôi như đang cầu cứu.

``\dots T-Tôi có nghe một lời đồn,'' Inari-san lẩm bẩm, giọng nhỏ nhưng đủ để tất cả nghe. ``Học sinh chuyển trường\dots\ phải được đối xử tốt. Nếu không\dots\ những điều tồi tệ sẽ xảy đến.''

\img{e1-2e}

Không gian như chìm xuống một tầng sâu hơn.

``Đó\dots\ chỉ là lời đồn thôi mà,'' cô bạn thuở nhỏ của cô cười gượng, nhưng ngay cả cô cũng không giấu nổi sự bồn chồn.

Tôi và Nii-chan liếc nhau. Câu nói ấy\dots\ quen thuộc đến rợn người.

``\dots Không lẽ\dots'' tôi mở lời.

``Ừ,'' anh gật nhẹ. ``Đúng là cái lời nguyền đó.''

Thứ mà chúng tôi đã nghe\dots\ trong chính giấc mơ của mình đêm trước.

Cùng với hình ảnh cô gái tóc nâu, mắt xanh, mặc bộ đồng phục đen.

Misora Shino.

\img{e1-2f}

Cô gái khiến mọi chuyện trở nên ma quái. Nhưng\dots\ chẳng phải mục tiêu của cô ta là Sora sao?

Vậy tại sao chúng tôi vẫn nhớ rõ ràng đến vậy?

Có lẽ\dots

\dots Không sớm thì muộn, chúng tôi sẽ còn gặp lại cô ta nữa.

\end{document}